\documentclass[11pt,letterpaper]{article}
\usepackage[margin=0.88in,headheight=15pt]{geometry}
\usepackage[T1]{fontenc}
\usepackage[utf8]{inputenc}
\usepackage{lmodern,microtype,amsmath,amssymb,amsthm,mathtools}
\usepackage{xcolor,booktabs,longtable,array,enumitem,graphicx}
\usepackage{fancyhdr,listings,xurl}
\usepackage[colorlinks=true,linkcolor=navy,citecolor=teal,urlcolor=teal]{hyperref}
\definecolor{navy}{HTML}{18344B}
\definecolor{teal}{HTML}{007C83}
\definecolor{pale}{HTML}{EDF5F5}
\definecolor{graytext}{HTML}{506171}
\hypersetup{pdftitle={Beyond Dirichlet GRH: transferable methods and concrete targets},
 pdfauthor={Research analysis prepared for Vladimir Reshetnikov}}
\pagestyle{fancy}\fancyhf{}
\lhead{\small\textcolor{graytext}{Beyond Dirichlet GRH}}
\rhead{\small\textcolor{graytext}{25 November 2026}}
\cfoot{\small\thepage}
\setlength{\parindent}{0pt}
\setlength{\parskip}{5.5pt}
\setlength{\emergencystretch}{3em}
\setcounter{tocdepth}{2}
\setlist[itemize]{topsep=3pt,itemsep=3pt,leftmargin=1.4em}
\setlist[enumerate]{topsep=3pt,itemsep=4pt,leftmargin=1.6em}
\newtheorem{theorem}{Theorem}[section]
\newtheorem{proposition}[theorem]{Proposition}
\newtheorem{lemma}[theorem]{Lemma}
\newtheorem{corollary}[theorem]{Corollary}
\theoremstyle{definition}\newtheorem{definition}[theorem]{Definition}
\theoremstyle{remark}\newtheorem{remark}[theorem]{Remark}
\newcommand{\Q}{\mathbb Q}\newcommand{\R}{\mathbb R}\newcommand{\C}{\mathbb C}
\newcommand{\Z}{\mathbb Z}\newcommand{\N}{\mathbb N}
\newcommand{\re}{\operatorname{Re}}\newcommand{\im}{\operatorname{Im}}
\newcommand{\Gal}{\operatorname{Gal}}\newcommand{\Cl}{\operatorname{Cl}}
\newcommand{\ord}{\operatorname{ord}}\newcommand{\Res}{\operatorname{Res}}
\newcommand{\one}{\mathbf 1}\newcommand{\code}[1]{\texttt{\small #1}}
\newcommand{\src}[1]{\nolinkurl{#1}}
\newcolumntype{P}[1]{>{\raggedright\arraybackslash}p{#1}}
\newcommand{\sha}{72fead102f663d5c7a407964a13dc0cf9d977207}
\newcommand{\keybox}[1]{\par\smallskip\noindent\colorbox{pale}{\parbox{\dimexpr\linewidth-2\fboxsep}{#1}}\par\smallskip}
\lstset{basicstyle=\ttfamily\footnotesize,breaklines=true,columns=fullflexible,
 frame=single,rulecolor=\color{graytext},showstringspaces=false,aboveskip=8pt,belowskip=8pt}

\begin{document}
\begin{center}
{\small\textcolor{teal}{REPOSITORY ANALYSIS AND RESEARCH PROGRAM}}\\[8pt]
{\LARGE\bfseries\textcolor{navy}{Beyond Dirichlet GRH}}\\[6pt]
{\large Transferable methods, established consequences,\\
and a concrete first Hecke target}\\[12pt]
{Prepared for Vladimir Reshetnikov}\quad\textbar\quad 25 November 2026\\[6pt]
{\small Repository: \url{https://github.com/lsernau/grh-lean}}\\
{\footnotesize Snapshot: \texttt{\sha}}
\end{center}

\section*{Executive assessment}
\addcontentsline{toc}{section}{Executive assessment}
\textbf{Yes: the result can resolve other questions by established implications, and its
internal method gives a concrete program for going further. These are different kinds of reuse.}
The clearest named consequence is Vinogradov's conjecture on the least quadratic nonresidue.
The repository supplies Dirichlet GRH; an existing theorem of Lamzouri, Li and Soundararajan
then gives $n(p)<(\log p)^2$ for every prime $p\ge5$. This implies
$n(p)\ll_\varepsilon p^\varepsilon$ with an explicit constant given below.
No new analytic proof of that implication is claimed here.

The most useful next target for the \emph{method} is the degree-two Hecke $L$-function attached
to a nontrivial order-three ideal-class character of $\Q(\sqrt{-23})$. Equivalently, after using
the inspected Dirichlet theorem, it is RH for the cubic field
$\Q(\alpha)$ with $\alpha^3-\alpha-1=0$. This target is outside the five exported statements.
We give a complete reduction to one precise one-sided principal-ideal prime-sum inequality.
The analytic reduction is proved in this report; the needed arithmetic inequality is not.

The proof architecture worth extracting is the combination of:
\begin{itemize}
\item many short integral vectors, with a quantitative bound on the dimension they span;
\item a lower bound for the same determinant, using different auxiliary twists at different primes;
\item cancellation of the common main term, leaving an error divided by the number of theta lifts;
\item a one-sided Mellin-transform criterion with a crucial sign condition on real residues.
\end{itemize}
The most important obstruction is equally specific: the current prime detector uses rational
residue classes and an auxiliary quadratic splitting character. It does not already distinguish
the order-three ideal classes required by the cubic example.

\keybox{\textbf{What this report establishes.} Exact scope of the selected release; an explicit
Vinogradov consequence; the abelian-number-field consequence; abstract transfer lemmas;
an exact reduction for the cubic target; and reproducible source, arithmetic and small Lean checks.
It does \emph{not} claim a new proof of Hecke GRH, Vandiver, Artin's full density conjecture,
or another conjecture beyond the stated consequences of Dirichlet GRH.}

\textbf{Verification basis.} The pinned source was inspected and its archive hash matched the
independent verification record. Sutherland's separate page for this exact short release reports
successful builds, statement comparisons and multiple kernel checks for both editions.
The full GRH project was not rebuilt in this workspace. Seven elementary transfer lemmas were
separately checked here with Lean~4.35.0 and the same pinned Mathlib; they use only the three
standard axioms. This distinction is maintained throughout.\cite{repo,shortchecks}

\clearpage
{\small\setlength{\parskip}{0pt}\tableofcontents}
\clearpage

\section{What the selected repository actually proves}
\subsection{Snapshot and evidence}
The analysis concerns the single commit \texttt{\sha}, dated 25 November 2026,
in Luke Sernau's \code{lsernau/grh-lean}. The repository describes itself as a shorter rewriting
of Levent Alp\"oge's argument developed with Claude. It contains independent Lake projects
\code{grh/} and \code{grh-compact/}. Both pin Lean~4.35.0 and Mathlib commit
\texttt{56a005855ea17e0b0958f73ec801ae140b0a1679}.\cite{repo,readme}

The exact source archive produced locally by \code{git archive HEAD} has SHA-256
\begin{center}\footnotesize\texttt{9138cf2f5862e64f823ed676e8d37333c0717e1a13cd8f74f6b735ab2d1ba57c}.
\end{center}
It agrees with the archive fingerprint published on Sutherland's short-release verification page.
That page was updated at 02:50 UTC on 26 November, which is 18:50 Pacific time on
25 November, the date used for this report. It reports successful sandboxed builds and
independent checking of the five theorem statements for \emph{this commit}, not only for the
earlier, much larger release. A stopped check on a smaller machine is identified there as a
memory-related interruption; the corresponding check completed on the larger machine.
These are independent reports consulted here, not local rerun results.\cite{shortchecks}

\begin{center}\small
\begin{tabular}{lrr}
\toprule
Local source audit & Readable edition & Compact edition\\
\midrule
Lean files under \code{src/} & 803 & 794\\
Physical source lines, including comments & 477,124 & 112,827\\
Approximate declarations, lexical count & 23,698 & 9,833\\
Flagged admissions/trust-extension tokens in code & 0 & 0\\
Local source modules reachable from \code{Main} & 803 & 794\\
\bottomrule
\end{tabular}
\end{center}
The audit removes nested comments and quoted strings before screening. It checks imports and
records per-file SHA-256 hashes. It is deliberately not a Lean parser or proof checker.
All modules being imported does not mean every declaration in them occurs in the final proof term.
The source itself notes that some declarations are unused.\cite{readme}

\subsection{The mathematical endpoints}
The crucial file is \src{grh/src/Statements.lean}. It imports only Mathlib. Its names refer to
Mathlib's actual Riemann zeta function, Dirichlet characters and analytically continued Dirichlet
$L$-functions. The five theorems in \src{grh/src/Main.lean} have the following meanings.

\begin{center}\small
\begin{tabular}{p{0.32\linewidth}p{0.62\linewidth}}
\toprule
Lean theorem & Mathematical assertion\\
\midrule
\src{riemannHypothesis} & Mathlib's standard RH statement.\\[3pt]
\src{riemannZeta_eq_zero_re} & A zero of $\zeta$ other than a negative even integer, with $s\ne1$, has real part $1/2$.\\[3pt]
\src{LFunction_ne_zero_of_half_lt_re} & For every $N\ge1$ and Dirichlet character $\chi\bmod N$, $L(s,\chi)\ne0$ if $\re s>1/2$, except for the principal character at $s=1$.\\[3pt]
\src{LFunction_eq_zero_re} & Every zero of every such $L$-function in $0<\re s<1$ lies on $\re s=1/2$.\\[3pt]
\src{LFunction_eq_zero_re_of_isPrimitive} & The same assertion restricted to primitive characters.\\
\bottomrule
\end{tabular}
\end{center}
The condition at $s=1$ handles the pole, where a totalized formal function needs a conventional
value. It is not an exemption for an ordinary zero. Imprimitive characters are covered, but
their additional finite Euler factors can have zeros on $\re s=0$; the strip statement is
therefore the appropriate endpoint.\cite{statements,main}

The pinned Mathlib was also inspected. Its zeta function agrees with the usual Dirichlet series
on $\re s>1$ and continues analytically away from~1. The function
\code{DirichletCharacter.LFunction} is defined by a Hurwitz-zeta continuation; it is not the
default value that a formal infinite sum may receive outside its convergence region.
The statement files of the two editions are byte-identical, and their final \code{Main.lean}
code is identical after ignoring comments and whitespace.

\textbf{Scope boundary.} Here ``GRH'' means \emph{Dirichlet} GRH. It does not, merely by its name,
assert RH for every Dedekind zeta function, every Hecke character over every number field,
or every automorphic $L$-function. Conversely, it is too dismissive to treat this as merely
a conjecture restated as an axiom: the inspected endpoint is unconditional, the source scan found
no such axiom, and the exact release has substantial independently reported kernel verification.

\subsection{What was and was not checked here}
The local work comprised source inspection, the lexical/import audit, comparison of the source
archive fingerprint, inspection of the pinned Mathlib definitions, exact finite arithmetic
checks, and a separate Lean check of the elementary lemmas supplied with this report.
The local Lean command for those lemmas exited successfully and printed only
\code{propext}, \code{Classical.choice} and \code{Quot.sound} for all seven declarations.

The full source-to-proof build and the large exported GRH proof were not rerun here.
An attempted download of the published export returned HTTP~403; it did not produce a local
proof-checker result. The report uses the published independent verification for that layer.
In particular, the small accompanying Lean file does not certify the deep theta-lift argument.
This is a mathematical transfer analysis with reproducible supporting checks, rather than
a claim to have personally read or rebuilt every proof in the development.

\section{Questions already settled by transferring the result}
\subsection{Vinogradov's least quadratic nonresidue conjecture}
For an odd prime $p$, let $n(p)$ be the smallest positive integer which is not congruent to a
nonzero square modulo~$p$. The familiar conjecture is
\begin{equation}\label{eq:vinogradov}
  \forall\varepsilon>0\quad n(p)\ll_\varepsilon p^\varepsilon.
\end{equation}
Tao's paper identifies this precise conjecture and recalls its implication from GRH.
Lamzouri--Li--Soundararajan, Corollary~1.1, give the explicit conditional bound
\begin{equation}\label{eq:lls}
  n(p)<(\log p)^2\qquad(p\ge5\text{ prime}).
\end{equation}
Their required Dirichlet-GRH hypothesis is supplied by the selected repository.
Thus the formerly conditional argument becomes a theorem once that verified result is used.
This is an immediate, substantive positive answer to the question about other conjectures;
the implication itself is classical, not a new discovery of this report.\cite{tao,lls}

\begin{corollary}[Explicit form of the Vinogradov consequence]\label{cor:vinogradov}
Using the repository's Dirichlet-GRH theorem and the Lamzouri--Li--Soundararajan bound, for every
$\varepsilon>0$ and every odd prime $p$,
\[
 n(p)\le C_\varepsilon p^\varepsilon,
 \qquad C_\varepsilon=
 \max\left\{2\,3^{-\varepsilon},\frac{4}{e^2\varepsilon^2}\right\}.
\]
\end{corollary}
\begin{proof}
For $u\ge0$, differentiating $u^2e^{-\varepsilon u}$ shows that its maximum is reached at
$u=2/\varepsilon$, with value $4/(e^2\varepsilon^2)$. Hence
$(\log p)^2\le4p^\varepsilon/(e^2\varepsilon^2)$. Apply \eqref{eq:lls} for $p\ge5$.
For $p=3$, the exact value $n(3)=2$ is covered by the other term in $C_\varepsilon$.
\end{proof}

This separates three pieces of evidence: the new GRH theorem, an existing published implication,
and the elementary maximization just proved. The accompanying script additionally checks the
small-prime range in \eqref{eq:lls}: all 428 primes $5\le p<3000$. It uses Euler's criterion and
exact rational lower bounds for $\log p$, so floating-point rounding is not a proof step.
That finite computation is useful corroboration; it does not establish the infinite assertion.

There are two different primitive-root questions that should not be confused with
\eqref{eq:vinogradov}. Bounding the \emph{least} primitive root as the prime modulus varies is
one problem. Proving a prescribed integer is a primitive root for a positive density of primes
is Artin's problem. The latter requires a different family of number-field $L$-functions in
Hooley's conditional argument; see Section~\ref{sec:limits}.

\subsection{Dedekind RH for every abelian extension of \texorpdfstring{$\Q$}{Q}}
\begin{theorem}[Abelian-number-field consequence]\label{thm:abelian}
For every finite abelian extension $K/\Q$, every zero of $\zeta_K(s)$ in $0<\re s<1$ has
real part~$1/2$, once the selected Dirichlet-GRH theorem is used.
\end{theorem}
\begin{proof}
By Kronecker--Weber and the standard character factorization,
\begin{equation}\label{eq:abelian}
 \zeta_K(s)=\prod_{\chi\in X_K}L(s,\chi),
\end{equation}
where $X_K$ is the finite family of primitive Dirichlet characters corresponding to $K$.
Each factor is holomorphic in the open critical strip. A zero of the product is a zero of at
least one factor. Dirichlet GRH gives its real part as $1/2$.\cite{kw}
\end{proof}

This includes quadratic and cyclotomic fields and all their subfields. The relevant word is
``abelian \emph{over $\Q$}.'' An extension abelian over a different number field can lead to
Hecke functions outside the selected theorem.

The primitive inducing characters in \eqref{eq:abelian} are important: replacing them indiscriminately
by imprimitive characters at a common modulus removes Euler factors and changes the exact identity.
The finite-product zero-transfer lemma in the accompanying Lean file formalizes the elementary
last step, with the analytic factorization explicitly assumed. It does not formalize
Kronecker--Weber or manufacture a missing Hecke factorization.

\subsection{Other useful consequences and their novelty status}
Two further precise consequences deserve a place in an applications project.
First, Lamzouri--Li--Soundararajan, Corollary~1.2, give for $q>3$ and $(a,q)=1$
\[
 P(a,q)\le\bigl(\varphi(q)\log q\bigr)^2,
\]
where $P(a,q)$ is the least prime congruent to $a\pmod q$. Second, the repository's
half-plane theorem directly eliminates real Dirichlet zeros in $(1/2,1)$; it does not assert
nonvanishing at the central point.\cite{lls,statements}

These consequences differ from proving Linnik's conjectural nearly linear bound in the modulus.
The displayed least-prime bound is roughly quadratic in $q$, with logarithmic factors.
The result is powerful but its exponent should not be silently improved.

Robin's and Lagarias's divisor-sum inequalities are also consequences of RH, but they are
especially poor candidates for a novelty claim here: the original verification project has
already combined its RH proof with an independent formalization and checked those arithmetic
statements, as well as a corrected cyclotomic ideal inequality.\cite{originalchecks}

The appropriate deliverable for these immediate consequences would be an independently
stated applications library: explicit hypotheses, standard mathematical definitions, and checked
bridges from the exported Dirichlet theorem. That is useful formalization work even when
the mathematical implication has been known for decades.

\section{The internal mechanism that may transfer}
\subsection{The arithmetic statistic}
The analytic end of the proof studies a smoothed prime-power remainder of the form
\begin{equation}\label{eq:primeR}
 R_{d,q}(x)=
 \sum_{\substack{n\le x\\n\equiv1\pmod q}}
 \Lambda(n)(1-\chi_d(n))\left(\frac1n-\frac1x\right)
 -\frac{\log x}{\varphi(q)}-c_{d,q}.
\end{equation}
Here $d$ is an auxiliary prime, $\chi_d$ is its quadratic character in the conventions of the
source, and $c_{d,q}$ is the normalizing constant. In the admissible setting $d\equiv7\pmod8$,
the corresponding quadratic field is $\Q(\sqrt{-d})$.
The coefficients are nonnegative. The smoothing factor vanishes at $n=x$ and yields the
denominator $s(s-1)$ in a Mellin transform; it is part of the proof mechanism, not cosmetic.
See \src{New/Reduction/Zeros/PrimeError.lean}.\cite{primeerror,criterion}

The geometric part first works with special moduli $q_0$ divisible by $18480$ and primes
$d\equiv q_0/2-1\pmod{q_0}$. The final reduction covers arbitrary Dirichlet characters.
The intermediate restrictions therefore should neither be mistaken for restrictions on the
final theorem nor ignored when trying to reuse the intermediate construction.

\subsection{Three estimates for one arithmetic lattice}
Let $D$ be the dimension of the target space, $n$ the number of independent forms selected,
$m$ the number of theta lifts, and $M$ the common main contribution. A finite-dimensional
arithmetic lattice comes with ordinary archimedean lengths and integrality conditions at finite
primes. Its squared covolume is a Gram determinant. The source compares that determinant
in two ways, while charging the denominators needed to make vectors integral.

\begin{enumerate}
\item \textbf{Short forms (Theorem A).} High derivatives of CM theta functions have
large prime divisibility. After the prescribed normalization this yields short archimedean
vectors and bounds on their denominators.
\item \textbf{Enough forms (Theorem B).} A second-moment argument gives translates spanning
at least $(1-4/b^2)D$ dimensions. A single nonzero lift would not provide this quantitative input.
\item \textbf{Integral forms are long (Theorem C).} Product-formula and determinant estimates
give a lower bound for the same covolume plus twice its index cost.
\end{enumerate}
These are connected in \src{New/Reduction/ShortAndEnough.lean},
\src{New/Determinant/Twist/LongForms.lean}, and
\src{New/Reduction/MainInequality.lean}.\cite{shortenough,longforms,mainineq}

Writing $V$ for a logarithmic squared-covolume bound and $\delta$ for an index bound,
the comparison has the schematic form
\begin{equation}\label{eq:sandwich}
 D(M-B)\ \le\ V+2\delta\ \le\ n(M-mR+A)+2DI,
 \qquad \theta D\le n\le D.
\end{equation}
All quantities refer to compatible objects. Cancellation of $M$ leaves a bound for $R$
with a gain of $m$. The rigorous algebraic form is proved in Section~\ref{sec:abstract}.

In the actual data the source obtains
\begin{equation}\label{eq:star}
 m\,\mathrm{totalRemainder}\le f a(\log b)^5,
\end{equation}
where $f$ counts infinite places and $a=xb$ is a derivative order. The calculation keeps track
of the fraction of the target spanned, the doubled index cost, the error in the long-form bound,
and the sign $M\ge0$. Each is needed for the cancellation.\cite{mainineq}

\subsection{Why the chain has square-root length}
After changing from the total remainder to \eqref{eq:primeR}, the source proves
\begin{equation}\label{eq:cutoffs}
 R_{d,q_0}(x)\le
 \frac{\varphi(q_0)(\log b)^6}{m}
 +|R_{d,q_0}(x_1)|+(\log b)|R_{d,1}(x_1)|.
\end{equation}
The admissible scales include $m\in[\sqrt{x},2\sqrt{x}]$. For a small fixed $\eta>0$ they can
be chosen so that $\log b\le2x^\eta$ and the two remote-cutoff errors are at most $x^{-2}$.
Consequently
\[
 R_{d,q_0}(x)\le(64\varphi(q_0)+3)x^{-1/2+6\eta}
\]
for sufficiently large integers $x$, followed by a controlled passage to real $x$.
This becomes an eventual one-sided $x^{-1/2+\varepsilon}$ bound.
The threshold depends on the fixed parameters; the construction does not by itself print
an explicit uniform onset for every modulus.\cite{onesided}

Alp\"oge's announcement explains the guiding analogy: each additional theta lift plays a role
like extra multiplicity in Stepanov's method over function fields. The square-root number
of lifts then supplies the desired square-root saving. That analogy is useful, but the
existence of compatible data and control of the errors are the substance of the argument.
One cannot increase $m$ arbitrarily and infer a stronger zero-free region.\cite{announcement}

\subsection{A different auxiliary twist at each prime}
The lower-bound interface is more flexible than choosing one globally favorable twist.
For each prime and each relevant embedding into an algebraic closure of a $p$-adic field,
it permits a suitable twist. Archimedean size is still measured using one fixed untwisted
determinant. Ratios between the twisted and untwisted determinants are bounded through
a height construction and congruence-denominator estimates.

The definition \code{HeightLe} includes a \emph{common nonzero multiplier} and bounds at all
archimedean conjugates and finite places. This is what makes the product formula usable.
One cannot assign unrelated local quantities at different primes and call the result a
global lattice. The compatibility is an explicit mathematical requirement. The generic
objects are in \src{New/Reduction/Lattices/Heights.lean} and
\src{New/Reduction/Lattices/TwistedLattice.lean}; existence of the required regular families
is proved for the specialized data in \src{New/Congruences/FitTwistsExist.lean}.
\cite{heights,twisted,regulartwists}

This suggests a reusable technique for arithmetic inequalities where no single twist makes
every congruence favorable. It does not yet give a theorem about arbitrary Selmer groups or
arbitrary $L$-functions. The determinant comparison and the arithmetic construction must
both survive an extension.

\section{Two transfer principles, with proofs}\label{sec:abstract}
\subsection{The quantitative determinant comparison}
\begin{proposition}[Determinant amplification]\label{prop:comparison}
Let $D,m,\theta>0$, $\theta D\le n\le D$, and $M,B,I\ge0$.
Let $A,R,V,\delta\in\R$, where $A$ and $\delta$ need not be nonnegative.
Suppose
\[
 V\le n(M-mR+A),\qquad \delta\le DI,\qquad
 D(M-B)\le V+2\delta.
\]
Then
\begin{equation}\label{eq:sharp}
 R\le\frac{A+(B+2I)/\theta}{m}.
\end{equation}
\end{proposition}
\begin{proof}
The hypotheses give
\[
 nmR\le(n-D)M+nA+D(B+2I)\le nA+D(B+2I),
\]
because $n\le D$ and $M\ge0$. Also $n\ge\theta D>0$, so division by $nm>0$ yields
\[
 R\le\frac{A+(D/n)(B+2I)}m
 \le\frac{A+(B+2I)/\theta}m.
\]
The final inequality uses $D/n\le1/\theta$ and $B+2I\ge0$.
\end{proof}
This proposition is independently formalized as
\code{GrhTransfer.sharp\_determinant\_sandwich} in the accompanying Lean file.
The statement makes the reusable part exact: the same main term, a positive fraction of
the target dimension, and error controlled relative to the number of lifts.
No theta-theoretic assertion is hidden in this lemma.

\begin{corollary}[Completion cost]
Suppose the last lower bound is replaced by
\[
 D(M-B_0)-(D-n)K\le V+2\delta,
\]
where $B_0,K\ge0$, and $n\ge(1-\eta)D$ for $0\le\eta<1$. Then
\[
 R\le\frac{A+(B_0+2I+\eta K)/(1-\eta)}m.
\]
\end{corollary}
\begin{proof}
Use $D-n\le\eta D$, so $B=B_0+\eta K$ is an admissible loss in
Proposition~\ref{prop:comparison}, and set $\theta=1-\eta$.
\end{proof}
Thus ``almost spanning'' is useful only together with an affordable cost for the missing
dimensions. Its rate cannot be discarded when transplanting the argument.

If the numerator of \eqref{eq:sharp} is $O_\varepsilon(X^\varepsilon)$ and
$m\ge cX^\mu$ with $0<\mu<1$, the conclusion is
$R(X)\le O_\varepsilon(X^{-\mu+\varepsilon})$.
With the analytic assumptions below, this is a route to a zero-free half-plane
$\re s>1-\mu$. Full square-root amplification corresponds to $\mu=1/2$.
This is a specification of what an extension must prove, not an assertion that those
inputs already exist for another family. The Lean file checks the real-power bookkeeping
with the amplification and error estimates left as explicit hypotheses.

\subsection{The one-sided analytic criterion}
\begin{proposition}[One-sided pole exclusion]\label{prop:landau}
Let $E:[0,\infty)\to\R$ be measurable, locally bounded and of at most exponential growth.
Suppose
\[
 \mathcal H(s)=\int_0^\infty E(v)e^{(1-s)v}\,dv
\]
initially converges in some right half-plane and has a meromorphic continuation to
$\re s>0$ with at most simple poles. Fix $\beta>0$. Assume
\begin{enumerate}
\item $E(v)\le C e^{(\beta-1)v}$ for all sufficiently large $v$;
\item the residue of $\mathcal H$ at every real point $u>\beta$ is nonnegative,
where the residue is zero at a regular point.
\end{enumerate}
Then $\mathcal H$ is holomorphic throughout $\re s>\beta$.
\end{proposition}
\begin{proof}
By enlarging a nonnegative constant $K$, local boundedness and the eventual estimate give
$K e^{(\beta-1)v}-E(v)\ge0$ for every $v\ge0$. Consequently
\[
 \mathcal F(s)=\frac K{s-\beta}-\mathcal H(s)
 =\int_0^\infty\bigl(K e^{\beta v}-e^v E(v)\bigr)e^{-sv}\,dv
\]
is a Laplace transform of a nonnegative function. Let $\sigma_c$ be its abscissa of
convergence; it is bounded above. If $\sigma_c>\beta$, Landau's theorem says that the
real point $\sigma_c$ is a singularity. Meromorphic continuation and the simple-pole
assumption make it a genuine simple pole. Since $\mathcal F(x)\ge0$ for real
$x>\sigma_c$, its residue there is positive. But the term $K/(s-\beta)$ is regular
there, so the residue is minus that of $\mathcal H$, hence nonpositive. This is impossible.
Thus $\sigma_c\le\beta$ (including the entire-transform possibility), and
$\mathcal F$, hence $\mathcal H$, is holomorphic for $\re s>\beta$.
\end{proof}
This is the classical Landau positivity mechanism in a transfer-friendly form, not a new
analytic theorem. The source proves an appropriate version in
\src{Lib/Analytic/Landau.lean}. Its application to the prime remainder is in
\src{New/Reduction/Zeros/Criterion.lean}.\cite{landau,criterion}

\subsection{Two hypotheses that cannot be omitted}
An upper bound alone does not exclude poles. For $0<\beta<\alpha<1$, take
\[
 E(v)=-e^{(\alpha-1)v}.
\]
Then $E(v)\le e^{(\beta-1)v}$ for every $v\ge0$, but
\[
 \mathcal H(s)=-\frac1{s-\alpha}
\]
has a pole to the right of $\beta$. Its real residue is negative, exactly the prohibited case.

There is a separate cancellation problem. A difference of logarithmic derivatives can fail
to see a zero shared by its two factors: if $F(s)=G(s)=s-\rho$, then
$F'/F-G'/G=0$, although both functions vanish at $\rho$.
This is why a theorem about a remainder must identify its residues and prevent cancellation,
rather than merely exhibit a meromorphic transform.

In the repository the residues have the shape
\begin{equation}\label{eq:resdifference}
 \frac{m(\rho)-m_d(\rho)}{\varphi(q)\rho(1-\rho)},
\end{equation}
where $m$ sums zero orders in an untwisted Dirichlet family and $m_d$ sums orders in its
quadratic twists. A zero-density argument selects an auxiliary $d$ for which the twisted
family has no zero in the finite region relevant to a putative zero $\rho$.
The real residues are then nonnegative, and the numerator at $\rho$ cannot cancel.
Proposition~\ref{prop:landau} forces it to be zero.\cite{criterion}

For this application, the geometric estimate and the auxiliary nonvanishing theorem are
separate ingredients. A transfer using a difference such as \eqref{eq:resdifference} needs
both. The direct ideal-class selector in the next section has a different residue structure
and avoids the auxiliary zero-avoidance step altogether.

\section{A concrete next target: the cubic field of discriminant \texorpdfstring{$-23$}{-23}}
\subsection{Why this is a useful first extension}
Set
\[
 F=\Q(\alpha),\quad\alpha^3-\alpha-1=0,\qquad
 K=\Q(\sqrt{-23}),\qquad H=KF.
\]
The class group of $K$ has order three. Let $\psi$ be either of its two nontrivial characters.
The target is the Hecke $L$-function $L_K(s,\psi)$, whose Euler product is over prime
\emph{ideals} of $K$, rather than rational residue classes. This is a degree-two
dihedral $L$-function over $\Q$.

There are exact factorizations
\begin{align}
 \zeta_F(s)&=\zeta(s)L_K(s,\psi),\label{eq:cubicfactor}\\
 \zeta_H(s)&=\zeta(s)L(s,\chi_{-23})L_K(s,\psi)^2.\label{eq:hilbertfactor}
\end{align}
Thus, after using the repository's RH and Dirichlet GRH, the three assertions
\begin{enumerate}
\item GRH for $L_K(s,\psi)$,
\item RH for $\zeta_F$,
\item RH for $\zeta_H$
\end{enumerate}
are equivalent in the open critical strip. None is among the five exported theorems.
The Hilbert-class-field identification is recorded in the classical literature; we give
the algebraic verification and the factorization argument explicitly.\cite{morishita}

\subsection{Algebraic verification}
The polynomial $x^3-x-1$ has no rational root and has discriminant
$-4(-1)^3-27(-1)^2=-23$. Therefore its normal closure has Galois group $S_3$, with
quadratic subfield $K$. The discriminant is squarefree, so $\Z[\alpha]$ is the full ring
of integers of~$F$.

The reduced positive definite binary quadratic forms of discriminant $-23$ are
\[
 (1,1,6),\qquad (2,1,3),\qquad (2,-1,3).
\]
Indeed the reduction bound is $3a^2\le23$, so $a=1$ or~2, and
$b^2-4ac=-23$ together with the reduction conditions gives exactly that list.
The usual form-class correspondence yields $h_K=3$.

Only the prime~23 can ramify in the normal closure, and
\[
 x^3-x-1\equiv(x-10)^2(x-3)\pmod{23}.
\]
The ramification is tame, with inertia a transposition in $S_3$.
Its intersection with $A_3=\Gal(H/K)$ is trivial, so $H/K$ is unramified at finite
places. There are no real places of~$K$. Hence $H/K$ is an unramified cyclic cubic
extension; because $h_K=3$, it is the Hilbert class field.

Let $\rho$ be the standard two-dimensional representation of $S_3$.
The permutation representation on the three embeddings of $F$ is $1\oplus\rho$, and
the regular representation is $1\oplus\mathrm{sgn}\oplus2\rho$.
Moreover $\rho=\operatorname{Ind}_{A_3}^{S_3}\psi$.
Artin formalism and induction of Euler products give \eqref{eq:cubicfactor} and
\eqref{eq:hilbertfactor}. No unproved general Artin holomorphy conjecture is being used:
this induced representation corresponds to the classical entire Hecke function of the
nontrivial ideal-class character. Complex conjugation of ideals preserves norms and inverts
classes, so $L_K(s,\psi^{-1})=L_K(s,\psi)$.

The following local table makes the distinction from a product of Dirichlet factors visible.
For an unramified rational prime put $T=p^{-s}$.
\begin{center}\small
\begin{tabular}{llll}
\toprule
Frobenius in $S_3$ & $\operatorname{tr}\rho$ & Denominator for $L(s,\rho)$ & Denominator for $\zeta_F$\\
\midrule
Identity & $2$ & $(1-T)^2$ & $(1-T)^3$\\
Transposition & $0$ & $1-T^2$ & $(1-T)(1-T^2)$\\
Three-cycle & $-1$ & $1+T+T^2$ & $1-T^3$\\
\bottomrule
\end{tabular}
\end{center}
The accompanying exact-arithmetic script checks these polynomial identities, the character
decompositions and the ramified factor at~23. It also computes the first hundred coefficients
of the explicit theta-series difference
\[
 \frac12\left(\sum_{u,v\in\Z}q^{u^2+uv+6v^2}
                 -\sum_{u,v\in\Z}q^{2u^2+uv+3v^2}\right).
\]
The tested prime coefficients agree with
$\#\{x\bmod p:x^3-x-1=0\}-1$ for all 24 unramified primes through~97.
This finite agreement is a consistency check, not a proof of an analytic continuation
or a zero-location statement.

\subsection{An exact positive selector}
For a nonzero integral ideal $\mathfrak a$ of~$K$, write $N\mathfrak a$ for its norm and
$[\mathfrak a]$ for its ideal class. Character orthogonality gives
\begin{equation}\label{eq:selector}
 \one_{[\mathfrak a]=1}
 =\frac{1+\psi(\mathfrak a)+\psi(\mathfrak a)^2}{3}.
\end{equation}
The left side is zero or one, despite the complex characters on the right.
Let $\Lambda_K(\mathfrak p^j)=\log N\mathfrak p$ for a prime ideal $\mathfrak p$ and
$j\ge1$, and let $\Lambda_K$ vanish on other ideals. Define
\[
 \mathcal A(s)=\sum_{\mathfrak a}
 \frac{\Lambda_K(\mathfrak a)\one_{[\mathfrak a]=1}}{(N\mathfrak a)^s}.
\]
Then, for $\re s>1$,
\begin{equation}\label{eq:logderivative}
 \mathcal A(s)=-\frac13\frac{\zeta_H'(s)}{\zeta_H(s)}.
\end{equation}

Here is a direct local proof, which also checks prime powers. If the class of a prime ideal
$\mathfrak p$ has order $f\in\{1,3\}$, then $\mathfrak p$ splits into $3/f$ primes of
norm $(N\mathfrak p)^f$ in~$H$. Dividing the logarithmic derivative by~3 gives
coefficient $\log N\mathfrak p$ at exactly the powers $\mathfrak p^j$ with $f\mid j$.
Those are exactly the principal powers. Since $H/K$ is unramified, no finite Euler factor
needs removal. The same argument works with $1/h$ for any unramified abelian extension
of degree~$h$, using its trivial Artin-symbol selector.

\subsection{A fully specified one-sided inequality}
Define the Euler--Kronecker constant of $H$ by
\[
 \zeta_H(s)=\frac{\kappa_H}{s-1}+\kappa_H\gamma_H+O(s-1).
\]
The exact target remainder is
\begin{equation}\label{eq:targetR}
 \boxed{\displaystyle
 \mathcal R(X)=
 \sum_{N\mathfrak a\le X}\Lambda_K(\mathfrak a)\one_{[\mathfrak a]=1}
 \left(\frac1{N\mathfrak a}-\frac1X\right)
 -\frac{\log X}{3}+\frac{1+\gamma_H}{3}.}
\end{equation}
The sign and constant are fixed by the Laurent expansion; an arbitrary additive constant
would destroy the desired decaying bound.

\begin{theorem}[Principal-class criterion for the cubic target]\label{thm:cubiccriterion}
If, for every $\varepsilon>0$, there are constants $C_\varepsilon$ and $X_\varepsilon$ such that
\begin{equation}\label{eq:missing}
 \mathcal R(X)\le C_\varepsilon X^{-1/2+\varepsilon}
 \qquad(X\ge X_\varepsilon),
\end{equation}
then RH holds for $\zeta_H$. Consequently GRH holds for $L_K(s,\psi)$ and RH holds for
$\zeta_F$, once the repository's Dirichlet theorem is supplied.
\end{theorem}
\begin{proof}
For $\re s>1$, termwise integration of the absolutely convergent prime-power sum yields
\begin{align}
 \mathcal M(s)
 &:=\int_1^\infty\mathcal R(X)X^{-s}\,dX\\
 &=\frac{\mathcal A(s)}{s(s-1)}
   -\frac1{3(s-1)^2}+\frac{1+\gamma_H}{3(s-1)}.\label{eq:targetM}
\end{align}
Indeed, a term of norm $n$ contributes
\[
 \int_n^\infty\left(\frac1n-\frac1X\right)X^{-s}\,dX
 =\frac{n^{-s}}{s(s-1)}.
\]
From \eqref{eq:logderivative},
$\mathcal A(s)=1/(3(s-1))-\gamma_H/3+O(s-1)$.
Substitution into \eqref{eq:targetM} cancels both poles at~1.
Elsewhere in $\re s>0$, its only possible poles occur at zeros $\rho$ of $\zeta_H$.
If the multiplicity is $m_H(\rho)$, the residue is
\begin{equation}\label{eq:targetres}
 \Res_{s=\rho}\mathcal M(s)
 =\frac{m_H(\rho)}{3\rho(1-\rho)}.
\end{equation}
These poles are simple even if the zeta zero is multiple. Every real residue with
$0<\rho<1$ is positive; there are no zeros on or to the right of $\re s=1$.

Set $E(v)=\mathcal R(e^v)$. It is locally bounded and of at most exponential growth.
For example, the norm-$n$ coefficient of the selected ideal von Mangoldt sum is bounded
by $2\Lambda(n)$, since $K$ is quadratic, so even a crude bound suffices.
The substitution $X=e^v$ identifies $\mathcal M(s)$ with the transform in
Proposition~\ref{prop:landau}. For $1/2<\beta<1$, apply \eqref{eq:missing} with
$\varepsilon=\beta-1/2$. The proposition makes $\mathcal M$ holomorphic in $\re s>\beta$.
Equation~\eqref{eq:targetres} then excludes every zeta zero there.
Let $\beta\downarrow1/2$ and use the Dedekind-zeta functional equation to exclude zeros
to the left of the line within the critical strip. The factorizations
\eqref{eq:cubicfactor}--\eqref{eq:hilbertfactor} give the remaining conclusions.
\end{proof}

\keybox{\textbf{The open work in this proposed transfer is \eqref{eq:missing}.}
The theorem above proves the analytic bridge. It does not prove the one-sided estimate,
and it is not presented as a new equivalent formulation unknown to the literature.
Its value is to give an exact mathematical target for extending the repository's construction.}

Unlike the difference-of-orders construction in \eqref{eq:resdifference}, this direct selector
has positive real residues automatically. It therefore requires no auxiliary zero-avoidance
theorem. If a geometric construction instead produces
$(1-\eta(\mathfrak a))\one_{[\mathfrak a]=1}$ for a quadratic Hecke character $\eta$,
the transform again contains a difference of zero multiplicities, and a separate
noncancellation argument becomes necessary. Those two routes should be kept distinct.

\subsection{The specific arithmetic obstruction}
Why not simply replace Dirichlet characters by ideal-class characters in the existing files?
The new information occurs at different primes. If a rational prime $p$ is inert in
$K=\Q(\sqrt{-23})$, the unique prime ideal above it is $(p)$ itself, which is principal,
and has norm $p^2$. Thus $\psi((p))=1$. The distinction between the two nontrivial class
characters and the trivial one comes from split prime ideals.

The current local divisibility construction is organized around primes where the auxiliary
quadratic character is not~1, especially inert primes, together with rational Frobenius actions
on cyclotomic embeddings. It does not already detect which nontrivial ideal class a split
prime of $K$ occupies. Moreover the existing global parameters require
$q_0\ge18480$ and $q_0\mid2(d+1)$, so setting $d=23$ is not even admissible.
See the data definitions and the divisibility-count theorem.\cite{fields,divisibility}

This is a limitation of the \emph{naive specialization}, not an impossibility theorem.
A new construction could use a separate auxiliary CM field, or an ideal-class-sensitive
family of CM periods. It would need to recover a nonnegative Gram construction, compatible
integral structures, the principal-class divisibility count and the quantitative error budget.
The proximity of CM theta theory makes this target plausible; none of those requirements
is discharged by the existence of a Hecke $L$-value somewhere inside the current norm formulas.

\section{A practical research program}
\subsection{Ranked tasks and stopping criteria}
\begin{longtable}{P{0.05\linewidth}P{0.26\linewidth}P{0.43\linewidth}P{0.16\linewidth}}
\toprule
 & Task & Concrete success criterion & Assessment\\
\midrule\endfirsthead
\toprule & Task & Concrete success criterion & Assessment\\\midrule\endhead
1 & Extract generic height, lattice and pole interfaces & Independent definitions, documented hypotheses, theorem-level dependency cuts, and exact comparison bounds & High value; bounded scope\\[5pt]
2 & Formalize direct applications & Vinogradov and least-prime bounds from the exported theorem and explicit published implications & Mathematically settled implication; formal work remains\\[5pt]
3 & Build the discriminant-$-23$ detector & A local divisibility theorem whose weighted count is the principal-class sum in \eqref{eq:targetR} & First serious new mathematical hurdle\\[5pt]
4 & Rebuild quantitative A/B/C for that detector & Almost-spanning forms, affordable completion, compatible denominator clearing and matching norm main terms & Substantial research\\[5pt]
5 & Prove \eqref{eq:missing} and assemble & Square-root amplification with losses $X^{o(1)}$, then Theorem~\ref{thm:cubiccriterion} & Would settle the chosen Hecke instance\\[5pt]
6 & Extend to fixed-field ray-class families & Control ramification and constants uniformly enough for the intended application & Later family-level goal\\
\bottomrule
\end{longtable}

The third task is the best early test of feasibility. If the required ideal-class-sensitive
divisibility does not materialize, polishing a huge higher-rank extension is premature.
Conversely, a successful local detector would justify serious work on spanning and lower bounds.

\subsection{What partial progress would count}
A first extension need not deliver $m\asymp\sqrt X$. If it proves a compatible comparison with
$m\gg X^\mu$ for some fixed $0<\mu<1/2$, losses $X^{o(1)}$, and the residue conditions of
Proposition~\ref{prop:landau}, it would give a zero-free region
$\re s>1-\mu$ for the target. Whether that would be new must be checked against current
results for the precise family. A correct weaker strip is not automatically novel.

A useful intermediate theorem could instead bound the denominator cost of a projection to a
theta-lift subspace independently of troublesome even Selmer classes. Such a result might
improve an arithmetic application without changing any zero-free region. Here too the
representation, local conditions, quantifiers and dependence on the twist must be retained.
An existence theorem for a favorable twist does not imply a bound for every member of a family.

The generic infrastructure is valuable even if the new detector fails. The source contains
specialized constructions for class field theory, Galois representations, theta correspondences,
comparison theory and trace formulas. Some are deliberately weaker than their standard full
versions. Their extraction should begin with precise semantic comparisons, rather than treating
every declaration under \code{Lib} as an interchangeable formalization of an entire theory.
The guide explicitly treats its literature classification as provisional and makes no novelty
claim based on the namespace \code{New}.\cite{guide}

\section{Questions not resolved by the identified transfers}\label{sec:limits}
\subsection{Artin's primitive-root density problem}
Hooley's conditional route uses Dedekind RH for the fields
\[
 \Q(\zeta_k,a^{1/k}),\qquad k\text{ squarefree},
\]
with estimates as $k$ varies. These fields are generally nonabelian over~$\Q$.
The inspected theorem therefore does not discharge those hypotheses by substitution.
Proving one cubic Hecke instance would still be far from controlling the full Kummer family.
This is a precise gap in this route, not a claimed theorem of logical independence between
Dirichlet GRH and Artin's conjecture. Separate contemporary progress on weaker infinitude
statements would not by itself establish the full density asymptotic.\cite{perucca}

\subsection{Class numbers versus class-group generators}
For an imaginary quadratic field, the class-number formula involves
$L(1,\chi_D)$, which is a Dirichlet function; GRH-derived class-number estimates are therefore
within reach. Bounds for very small prime-ideal generators of a class group normally use
nontrivial class-character Hecke functions. These are not all products of rational Dirichlet
functions. The order-three character of $\Q(\sqrt{-23})$ makes the distinction concrete.
For order-two genus characters additional factorizations can occur, which is one reason
to choose order three as the first diagnostic example.

\subsection{Elliott--Halberstam, Goldbach and twin primes}
The usual Dirichlet-GRH error estimate has the shape
\[
 \psi(x;q,a)-\frac{x}{\varphi(q)}
 =O\!\left(\sqrt x\log^2(qx)\right).
\]
Summing this bound over $q\le Q$ yields only
$O(Q\sqrt x\log^2(xQ))$. It cannot by this operation provide the nearly $x$-sized modulus
ranges of Elliott--Halberstam. That requires additional average cancellation.
Likewise the construction controls single prime sums, not the two-prime correlations needed
for binary Goldbach or twin-prime asymptotics. This identifies what is absent; it does not
assert that no imaginative use of GRH could ever contribute to those problems.
Tao's account makes the distinction between the distributional hypothesis and the least
nonresidue consequence especially useful.\cite{tao}

\subsection{Vandiver, algebraic \texorpdfstring{$K$}{K}-groups, and BSD}
The original announcement mentions a weakened Vandiver-type input and speculative
$K$-group consequences. The selected rewriting must be read on its own terms.
In \src{New/Congruences/FirstOrder/DenominatorBound.lean}, the explanation of the denominator
bound expressly says that the even classes, where Vandiver's conjecture lives, do not enter.
The cyclotomic bound \code{exists\_length\_shaOdd\_le} is an upper bound for particular
odd cocycle modules, with specified parity and local conditions; it is not a vanishing theorem
for all even class-group components.\cite{announcement,denominators,oddsha}

This distinction is mathematically decisive. A bound on a module's length by a valuation of
a large integer does not force the module to vanish at primes dividing that integer. Nor does
finiteness imply cyclicity. The rewrite's success in \emph{avoiding} certain classes is not a
proof that those classes are zero. Claims about the original release or additional
$K$-theory arguments need a separate audit of their exact inputs.

The same caution applies to BSD. The proof constructs favorable twists with small Selmer
groups for a particular denominator budget. It does not equate the analytic rank of an arbitrary
elliptic curve with its Mordell--Weil rank, determine a regulator, or establish the full leading-term
formula. Zero localization also allows central zeros and multiple zeros. Simplicity of zeta
zeros and central nonvanishing need an additional mechanism not supplied by the five endpoints.

\section{Deliverables, reproducibility and precise claim boundaries}
The archive contains the complete \LaTeX{} source and compiled report, a pinned small Lean
project, exact-arithmetic scripts and their JSON outputs, the static source-audit script and
its per-file hashes, selected unchanged source files for inspection, upstream licensing notices,
and instructions for reproducing the checks.

\begin{center}\small
\begin{tabular}{p{0.31\linewidth}p{0.61\linewidth}}
\toprule
Artifact & What it establishes\\
\midrule
\src{scripts/static_audit.py} & Repeatable source-token and import checks for the pinned repository; no kernel verification.\\[4pt]
\src{scripts/nonresidue_checks.py} & Exact finite verification of $n(p)<\log^2p$ for all 428 primes $5\le p<3000$.\\[4pt]
\src{scripts/arithmetic_checks.py} & Exact finite algebra behind the discriminant-$-23$ example, Euler factors and tested theta coefficients.\\[4pt]
\src{lean/Transfer.lean} & Seven elementary comparison, scaling and finite-factorization lemmas checked in Lean; all deep arithmetic hypotheses explicit.\\[4pt]
\src{evidence/lean_transfer_check.log} & The local Lean check output and printed axiom dependencies.\\[4pt]
\src{evidence/provenance.json} & Source commit, archive fingerprint, pinned versions and verification scope.\\
\bottomrule
\end{tabular}
\end{center}

The accompanying proofs in Sections~\ref{sec:abstract} and~5 are human-readable mathematical
arguments. The full Landau argument and the Hecke-target criterion have not been formalized
in the small Lean supplement. No experiment on finitely many primes or zeros is being used
as proof of an infinite zero-location claim.

For a research project built around this material, the recommended first mathematical goal is
the principal-class detector for discriminant $-23$, followed by its quantitative A/B/C
estimates. For an immediate useful applications project, formalizing the Vinogradov and
abelian-Dedekind bridges is the natural starting point. The report supplies both a firm
positive consequence and a testable extension program, while locating the unsolved step
in a single explicit inequality.

\appendix
\section{Source map}
All paths in this table are relative to \src{grh/src/} at the pinned commit. The table is a
reading map, not a claim that the short listed files contain all transitive dependencies.
\begin{longtable}{P{0.54\linewidth}P{0.40\linewidth}}
\toprule Source path & Reason to read it\\\midrule\endfirsthead
\toprule Source path & Reason to read it\\\midrule\endhead
\src{Statements.lean}, \src{Main.lean} & Exact mathematical endpoints.\\[4pt]
\src{New/Reduction/MainInequality.lean} & Cancellation of main terms and complete error budget.\\[4pt]
\src{New/Reduction/ShortAndEnough.lean} & Short forms, index control, almost-spanning statement.\\[4pt]
\src{New/Determinant/Twist/LongForms.lean} & Lower covolume/index bound.\\[4pt]
\src{New/Reduction/Lattices/Heights.lean} & Common-multiplier height formalism.\\[4pt]
\src{New/Reduction/Lattices/TwistedLattice.lean} & Prime-dependent twists with a fixed determinant.\\[4pt]
\src{New/Congruences/FitTwistsExist.lean} & Existence of twists with controlled denominators.\\[4pt]
\src{New/Reduction/OneSidedBound.lean} & Square-root scales and the final prime-remainder estimate.\\[4pt]
\src{New/Reduction/Zeros/PrimeError.lean} & Arithmetic sum, Mellin transform and zero residues.\\[4pt]
\src{New/Reduction/Zeros/Criterion.lean} & Auxiliary zero avoidance plus Landau.\\[4pt]
\src{Lib/Analytic/Landau.lean} & General analytic sign criterion.\\[4pt]
\src{New/Reduction/Data/FieldsAndPlaces.lean} & Specialized quadratic/cyclotomic setting.\\[4pt]
\src{New/Derivatives/Divisibility/Count.lean} & The prime detector that would need extension.\\[4pt]
\src{New/Congruences/FirstOrder/DenominatorBound.lean} & Why the rewritten denominator bound avoids Vandiver's even classes.\\[4pt]
\src{Lib/Cyclotomic/OddSha.lean} & Exact odd-component bound and its limitations.\\
\bottomrule
\end{longtable}

\section{Reproduction commands}
The scripts use the Python standard library. The Lean supplement uses the official toolchain
and the pinned Mathlib specified in its Lake files.
\begin{lstlisting}
python3 scripts/nonresidue_checks.py > data/nonresidue_checks.json
python3 scripts/arithmetic_checks.py --output data/arithmetic_checks.json

git clone https://github.com/lsernau/grh-lean.git upstream
git -C upstream checkout 72fead102f663d5c7a407964a13dc0cf9d977207
python3 scripts/static_audit.py upstream > static_audit.json
git -C upstream archive HEAD | sha256sum

cd lean
lake exe cache get Mathlib.Tactic.Linarith Mathlib.Tactic.Ring \
  Mathlib.Analysis.SpecialFunctions.Pow.Real
lake env lean Transfer.lean
\end{lstlisting}
To rebuild the \emph{upstream} project, follow its own README in either edition:
\code{lake exe cache get}, \code{lake build}, and
\code{lake env lean scripts/PrintAxioms.lean}. Resource requirements and the more extensive
external-checker commands are documented in the upstream verification record. The archive's
default reproduction script does not silently launch that much larger task.

\clearpage
\begin{thebibliography}{99}\small
\bibitem{repo} L. Sernau et al., \emph{grh-lean}, selected source snapshot,
commit \texttt{72fead10}, 25 November 2026.
\url{https://github.com/lsernau/grh-lean/tree/72fead102f663d5c7a407964a13dc0cf9d977207}.
\bibitem{shortchecks} A. Sutherland, \emph{Verification files for the shorter Lean proof of RH and GRH},
updated 26 November 2026, 02:50 UTC; consulted 25 November Pacific time.
\url{https://math.mit.edu/~drew/shortgrh/}.
\bibitem{originalchecks} A. Sutherland, \emph{Verification files for the Lean proof of the Riemann Hypothesis},
including independently checked arithmetic consequences.
\url{https://math.mit.edu/~drew/grh/}.
\bibitem{readme} Selected repository, readable-edition README.
\url{https://github.com/lsernau/grh-lean/blob/72fead102f663d5c7a407964a13dc0cf9d977207/grh/README.md}.
\bibitem{guide} Selected repository, reading guide.
\url{https://github.com/lsernau/grh-lean/blob/72fead102f663d5c7a407964a13dc0cf9d977207/grh/GUIDE.md}.
\bibitem{statements} Selected repository, \emph{Statements.lean}.
\url{https://github.com/lsernau/grh-lean/blob/72fead102f663d5c7a407964a13dc0cf9d977207/grh/src/Statements.lean}.
\bibitem{main} Selected repository, \emph{Main.lean}.
\url{https://github.com/lsernau/grh-lean/blob/72fead102f663d5c7a407964a13dc0cf9d977207/grh/src/Main.lean}.
\bibitem{announcement} L. Alp\"oge, announcement and remarks on the RH formalization, November 2026.
\url{https://alpo.ge/remark.pdf}.
\bibitem{lls} Y. Lamzouri, X. Li and K. Soundararajan, \emph{Conditional bounds for the least
quadratic non-residue and related problems}, Mathematics of Computation 84 (2015), 2391--2412;
arXiv version 3, 8 December 2016. Corollaries 1.1 and 1.2.
\url{https://arxiv.org/pdf/1309.3595}.
\bibitem{tao} T. Tao, \emph{The Elliott--Halberstam conjecture implies the Vinogradov least quadratic
nonresidue conjecture}, Algebra \& Number Theory 9 (2015), 1005--1034.
\url{https://msp.org/ant/2015/9-4/ant-v9-n4-p09-s.pdf}.
\bibitem{kw} A. Sutherland, \emph{The Kronecker--Weber theorem}, MIT 18.785, Lecture 20,
17 November 2021, including the primitive-character zeta factorization.
\url{https://math.mit.edu/classes/18.785/2021fa/LectureNotes20.pdf}.
\bibitem{morishita} M. Morishita, \emph{Integral representations of unramified Galois groups and matrix
divisors over number fields}, Osaka Journal of Mathematics 32 (1995), 565--575; example on p.~573.
\url{https://ir.library.osaka-u.ac.jp/repo/ouka/all/10928/ojm32_03_03.pdf}.
\bibitem{perucca} A. Perucca, \emph{A closed formula for the density in Artin's primitive root
conjecture over number fields}, discussion of Hooley's conditional theorem and its Kummer-field
hypotheses. \url{https://orbilu.uni.lu/bitstream/10993/66409/1/Perucca-Artin_formula.pdf}.
\bibitem{mainineq} Selected repository, \emph{The main inequality}.
\url{https://github.com/lsernau/grh-lean/blob/72fead102f663d5c7a407964a13dc0cf9d977207/grh/src/New/Reduction/MainInequality.lean}.
\bibitem{shortenough} Selected repository, \emph{Short and enough forms}.
\url{https://github.com/lsernau/grh-lean/blob/72fead102f663d5c7a407964a13dc0cf9d977207/grh/src/New/Reduction/ShortAndEnough.lean}.
\bibitem{longforms} Selected repository, \emph{Long integral forms}.
\url{https://github.com/lsernau/grh-lean/blob/72fead102f663d5c7a407964a13dc0cf9d977207/grh/src/New/Determinant/Twist/LongForms.lean}.
\bibitem{heights} Selected repository, \emph{Heights}.
\url{https://github.com/lsernau/grh-lean/blob/72fead102f663d5c7a407964a13dc0cf9d977207/grh/src/New/Reduction/Lattices/Heights.lean}.
\bibitem{twisted} Selected repository, \emph{Twisted lattices}.
\url{https://github.com/lsernau/grh-lean/blob/72fead102f663d5c7a407964a13dc0cf9d977207/grh/src/New/Reduction/Lattices/TwistedLattice.lean}.
\bibitem{regulartwists} Selected repository, \emph{Existence of regular twists}.
\url{https://github.com/lsernau/grh-lean/blob/72fead102f663d5c7a407964a13dc0cf9d977207/grh/src/New/Congruences/FitTwistsExist.lean}.
\bibitem{onesided} Selected repository, \emph{One-sided bound}.
\url{https://github.com/lsernau/grh-lean/blob/72fead102f663d5c7a407964a13dc0cf9d977207/grh/src/New/Reduction/OneSidedBound.lean}.
\bibitem{primeerror} Selected repository, \emph{Prime error}.
\url{https://github.com/lsernau/grh-lean/blob/72fead102f663d5c7a407964a13dc0cf9d977207/grh/src/New/Reduction/Zeros/PrimeError.lean}.
\bibitem{criterion} Selected repository, \emph{One-sided zero-free criterion}.
\url{https://github.com/lsernau/grh-lean/blob/72fead102f663d5c7a407964a13dc0cf9d977207/grh/src/New/Reduction/Zeros/Criterion.lean}.
\bibitem{landau} Selected repository, \emph{Landau's analytic criterion}, implementing the
classical positivity argument.
\url{https://github.com/lsernau/grh-lean/blob/72fead102f663d5c7a407964a13dc0cf9d977207/grh/src/Lib/Analytic/Landau.lean}.
\bibitem{fields} Selected repository, \emph{Fields and places}.
\url{https://github.com/lsernau/grh-lean/blob/72fead102f663d5c7a407964a13dc0cf9d977207/grh/src/New/Reduction/Data/FieldsAndPlaces.lean}.
\bibitem{divisibility} Selected repository, \emph{Divisibility count}.
\url{https://github.com/lsernau/grh-lean/blob/72fead102f663d5c7a407964a13dc0cf9d977207/grh/src/New/Derivatives/Divisibility/Count.lean}.
\bibitem{denominators} Selected repository, \emph{Denominator bound}.
\url{https://github.com/lsernau/grh-lean/blob/72fead102f663d5c7a407964a13dc0cf9d977207/grh/src/New/Congruences/FirstOrder/DenominatorBound.lean}.
\bibitem{oddsha} Selected repository, \emph{Odd cocycle bounds}.
\url{https://github.com/lsernau/grh-lean/blob/72fead102f663d5c7a407964a13dc0cf9d977207/grh/src/Lib/Cyclotomic/OddSha.lean}.
\end{thebibliography}
\end{document}
